\chapter*{\large ВВЕДЕНИЕ}  
\addcontentsline{toc}{chapter}{ВВЕДЕНИЕ}

Генерация процедурных изображений и карт является одной из важных задач компьютерной графики, моделирования виртуальных миров и разработки игровых приложений. При создании вымышленных ландшафтов, островов, материков, пещер и других природных структур требуется получать изображения, которые выглядят случайными, но при этом сохраняют правдоподобную форму и внутреннюю закономерность. Ручное построение таких объектов занимает много времени, поэтому на практике широко используются алгоритмы процедурной генерации.

Особое место среди таких методов занимают алгоритмы шума и дискретные модели. Шум Перлина, предложенный Кеном Перлином в 1982 году, позволяет получать плавные псевдослучайные значения и часто применяется для генерации текстур, облаков, рельефа и карт высот. Позднее был разработан симплексный шум, который решает часть проблем классического шума Перлина и лучше масштабируется при переходе к пространствам большей размерности. Алгоритм Diamond-Square применяется для построения фрактальных карт высот, а клеточные автоматы позволяют моделировать развитие дискретных структур, например пещер или областей суши и воды.

Актуальность работы связана с тем, что перечисленные алгоритмы используются в компьютерной графике, генерации игровых уровней, моделировании природных процессов и создании процедурного контента. Их изучение позволяет понять, как из простых математических правил можно получать сложные визуальные структуры.

Целью курсовой работы является изучение алгоритмов генерации псевдослучайных структур и реализация приложения с графическим интерфейсом для создания, отображения и комбинирования различных видов шума.

Для достижения поставленной цели необходимо решить следующие задачи:
\begin{enumerate}
    \item рассмотреть теоретические основы шума Перлина, симплексного шума, алгоритма Diamond-Square и клеточных автоматов;
    \item реализовать указанные алгоритмы на языке Python;
    \item разработать графическое приложение для запуска алгоритмов и отображения результата;
    \item добавить возможность комбинирования нескольких шумовых карт, применения масок и бинаризации изображения;
    \item провести анализ полученных результатов и определить особенности каждого метода.
\end{enumerate}

Объектом исследования являются алгоритмы процедурной генерации изображений. Предметом исследования являются методы построения двумерных матриц значений, которые могут интерпретироваться как карты высот, карты суши и воды или текстуры.

Практическим результатом работы является приложение на Python с графическим интерфейсом, позволяющее генерировать шум Перлина, симплексный шум, фрактальную карту Diamond-Square, применять клеточные автоматы, комбинировать результаты и визуализировать их в виде полутонового изображения.
